\documentclass[14pt]{mmcs_article}
\usepackage[russian]{babel}
\usepackage{amsmath, amsthm, amsfonts, amssymb}
\usepackage{enumitem}
%\graphicspath{{images/}}%путь к рисункам

\begin{document}

% Титульные листы
% раскомментировать требуемое
%\include{Titul3k} % для курсовой
\include{TitulBak}% для работы бакалавра
%\include{TitulMag}% для работы магистра

\renewcommand{\contentsname}{Оглавление}

\tableofcontents

%=======================
\newpage
\addcontentsline{toc}{section}{Постановка задачи}

\section*{Постановка задачи}


В постановке задачи коротко (по пунктам) указывается, что необходимо сделать в рамках работы. Раздел <<Постановка задачи>> должен соответствовать заданию на курсовую или выпускную квалификационную работу, подписанному научным руководителем.


%=======================
\newpage
\addcontentsline{toc}{section}{Введение}
\section*{Введение}

На данный момент не существует точной модели, описывающей физические процессы перед землетрясением (Родкин, 2011). Тем не менее, исследования продолжаются, так как накоплено много надежных данных о различных признаках приближающегося землетрясения: изменениях электромагнитных полей, химического состава вод и газов, деформациях поверхности и других (Cicerone et al., 2009; Geller, 2007; Ma et al., 2021; Chiodini et al., 2020).

По-видимому, некоторые из этих признаков способны улавливать животные. Известно множество случаев, когда незадолго до подземных толчков наблюдалось необычное поведение: собаки начинали громко лаять, кошки и собаки пытались убежать, крысы покидали норы, пчелы — ульи, а птицы издавали частые крики (Kirschvink, 2000). Также зафиксированы случаи, когда жабы покидали места своего обитания до завершения серии повторных толчков (Grant et al., 2011). В некоторых исследованиях даже численно показана связь между активностью мышей за день до землетрясения и изменениями геомагнитного поля (Yokoi et al., 2003).

С точки зрения эволюции такая чувствительность могла закрепиться, потому что одно крупное землетрясение (даже если оно случается раз в 50 поколений) способно уничтожить 10% и более популяции. Моделирование методом Монте-Карло показало высокую вероятность того, что способность чувствовать приближение катастрофы закреплялась в ходе естественного отбора (Kirschvink, 2000).

Известно, что реакция животных на предвестники землетрясений связана с состоянием стресса или повышенной тревоги. Стресс напрямую влияет на цикл сна и бодрствования, а также на общую активность (Rampin et al., 1991; Meerlo et al., 1997; Palma et al., 2000; Adrien et al., 1991). Согласно теории Ганса Селье (Selye, 1950), стресс — это общая реакция организма на вредные воздействия. Ее механизм связан с выбросом определенных гормонов (кортикотропин-рилизинг-гормона, а затем глюкокортикоидов). Выделение этих гормонов также подчиняется суточным ритмам, что задает циклический характер активности у грызунов (Herman et al., 2016; Kalsbeek et al., 2012). Колебания уровня этих гормонов меняют структуру сна у крыс (Bradbury et al., 1998). При этом разные виды стресса приводят к разным изменениям в структуре сна — например, одни увеличивают долю быстрого сна, а другие — долю медленного (Rampin et al., 1991; Meerlo et al., 1997; Palma et al., 2000; Adrien et al., 1991; Sanford et al., 2010).

Сон млекопитающих достаточно хорошо изучен. В нем выделяют две основные фазы: быстрый сон (с быстрыми движениями глаз, REM) и медленный сон (Rytkonen, 2011; Saevskiy, 2025; Kostin, 2019; Vyazovskiy, 2005; Lyamin, 2008). Для быстрого сна характерны низкоамплитудные высокочастотные колебания на электроэнцефалограмме и расслабление мышц. Эти две фазы циклически сменяют друг друга, причем медленный сон всегда предшествует быстрому. У грызунов сон, как правило, состоит из множества повторяющихся эпизодов в течение суток, причем спят они в основном в светлое время суток (Franken et al., 1999; Weber, Dan, 2016).

Исходя из вышесказанного, в данной работе мы исследовали изменения сна крыс как возможный результат стресса, вызванного воздействием факторов, предшествующих землетрясениям. Для исследования был выбран быстрый сон как наиболее чувствительный показатель стресса животного. Отдельное внимание уделено тому, за какой промежуток времени до землетрясения наблюдаются эти изменения.

Применение методов машинного обучения к описанной задаче выглядит естественным шагом: мы имеем многомерный временной ряд физиологических показателей, предположительно связанный с редкими внешними событиями, и хотим обнаружить эту связь, не имея строгой физической модели процесса. Однако специфика предметной области накладывает ряд серьезных ограничений, которые необходимо учитывать при выборе и оценке моделей.

\textbf{Проблема сверхмалой выборки.} Землетрясения достаточной силы, ощутимые на территории расположения лаборатории, являются редкими событиями. Даже при длительном непрерывном мониторинге число таких событий, попадающих в период регистрации, обычно исчисляется единицами или первыми десятками. Это фундаментальное ограничение: в отличие от типичных задач машинного обучения, мы не можем рассчитывать на тысячи или даже сотни примеров целевого события. Большинство мощных моделей (глубокие нейронные сети, градиентный бустинг) требуют значительно большего объема данных для обучения без переобучения. В данной ситуации предпочтительными оказываются экономные по числу параметров методы: линейные модели с регуляризацией, байесовские модели с сильными априорными предположениями, метод опорных векторов для обнаружения выбросов. Байесовский подход, кроме того, позволяет естественным образом количественно оценить неопределенность, связанную с малым объемом данных, — вместо точечных прогнозов модель выдает распределение возможных значений параметров, ширина которого отражает степень нашей неуверенности.

\textbf{Неоднородность данных.} Обучающая и тестовая выборки в данной задаче заведомо не являются независимыми и одинаково распределенными (i.i.d.). Данные представляют собой непрерывный временной ряд, где соседние наблюдения сильно скоррелированы. Кроме того, сами условия сбора данных могут меняться: сезонные факторы влияют на поведение животных, старение особи меняет ее физиологию, возможны дрейфы регистрирующей аппаратуры. Наконец, разные землетрясения различаются по магнитуде, расстоянию до эпицентра, глубине очага и геологическому контексту, а значит, и по характеру предвестников, достигающих лаборатории. Все это означает, что модель, хорошо работающая на одном временном отрезке, может оказаться несостоятельной на другом, если не учитывать перечисленные источники вариативности. Иерархические байесовские модели частично решают эту проблему, позволяя включать в структуру модели параметры, описывающие индивидуальные особенности животных и характеристики конкретных событий.

\textbf{Специфика работы с временными рядами.} Поскольку данные являются временным рядом, прямое применение стандартных методов классификации или регрессии, предполагающих независимость наблюдений, приводит к некорректным выводам. Временная структура порождает автокорреляцию: если крыса сегодня показала аномально высокую долю быстрого сна, то и завтрашнее значение, вероятно, будет отклоняться от среднего просто в силу инерционности физиологических процессов. Неучет этой зависимости ведет к ложной оценке значимости найденных закономерностей. При кросс-валидации временных рядов нельзя использовать случайное разбиение на фолды — необходимо применять схемы с сохранением временного порядка (time series split, скользящее окно), чтобы обучение всегда происходило на прошлом, а тестирование на будущем. Альтернативно, в байесовском подходе можно явно моделировать автокорреляционную структуру остатков, например, с помощью авторегрессионных компонент или гауссовских процессов, что позволяет корректно отделить эффект предвестников от естественной временной зависимости.

\textbf{Проблема множественного тестирования.} В ситуации, когда данных мало, а потенциальных предикторов (различные геофизические параметры, их лаги, частотные компоненты) много, возникает риск ложноположительных находок. Перебор большого числа гипотез без учета поправки на множественность почти гарантирует обнаружение статистически значимых, но содержательно пустых корреляций. Байесовский подход смягчает эту проблему, поскольку апостериорное распределение естественным образом штрафует сложные модели через маргинальное правдоподобие, а использование регуляризующих априорных распределений (например, horseshoe prior или spike-and-slab prior) позволяет автоматически отсеивать неинформативные предикторы.

\textbf{Интерпретируемость и научная ценность.} В отличие от многих прикладных задач машинного обучения, где главный критерий — точность предсказания, в данной работе не менее важна интерпретируемость результата. Нас интересует не только сам факт обнаружения аномалии перед землетрясением, но и количественная оценка временного окна, в котором эта аномалия проявляется, направленность изменений и степень индивидуальной вариабельности. Байесовский вывод здесь предпочтительнее многих «черных ящиков», так как дает прямые вероятностные утверждения об интересующих величинах.
Таким образом, с учетом перечисленных ограничений, в данной работе мы фокусируемся на байесовских методах анализа временных рядов как на инструменте, позволяющем совместить обнаружение аномальных паттернов сна с корректным учетом неопределенности, порожденной малым объемом данных и их сложной временной структурой. Основное внимание уделяется не столько построению предсказательной модели землетрясений, сколько проверке гипотезы о наличии устойчивого физиологического отклика и характеризации его временных свойств.



%=======================
\newpage
\section{Описание полученных результатов}\label{dsfs}

\subsection{Материалы и методы}
\subsubsection*{Экспериментальные животные и условия содержания}

Эксперименты проводились на базе Камчатского филиала Единой геофизической службы РАН (г.~Петропавловск-Камчатский, станция «Институт»). Объектом исследования служили взрослые самцы крыс линии пасюк массой от 200 до 250~г. Животные содержались поодиночке в прозрачных пластиковых боксах размером $30 \times 30 \times 40$~см с перфорированными стенками. В помещении поддерживалась постоянная температура ($23 \pm 1^{\circ}$C) и фиксированный 12-часовой цикл освещения (12~ч свет / 12~ч темнота). Доступ к пище и воде был свободным. Все процедуры были одобрены комиссией по биоэтике Южного федерального университета.

\subsubsection*{Регистрация биоэлектрической активности}

Для оценки состояний сна и бодрствования непрерывно регистрировалась электрическая активность коры больших полушарий (ЭКоГ) и гиппокампа (ЭСКоГ). Для этого животным под наркозом вживляли электроды в область сенсомоторной коры и в поле СА1 гиппокампа. Электроды соединялись с разъемом, который фиксировался на поверхности черепа. Регистрацию начинали не ранее чем через неделю после операции, чтобы исключить влияние восстановительного периода.

Запись сигналов велась с помощью 32-канального усилителя PBX2/32sp-r/32fp300 (Plexon Corp., США) и аналого-цифрового преобразователя L-ACRD (модель E~14-440). Частота оцифровки составляла 1~кГц по каждому каналу.

\subsubsection*{Алгоритм классификации состояний сна и бодрствования}

Для извлечения информации о функциональном состоянии животных из зарегистрированных сигналов использовался автоматический алгоритм стадирования, реализованный на языке Python. Алгоритм является упрощенной версией метода, описанного в работе \cite{Saevskiy2025}(Saevski), и относит каждую эпоху записи к одному из трех состояний: бодрствование, медленный сон (МС) или парадоксальный сон (ПС). Процедура обработки состоит из следующих шагов:

\begin{enumerate}
	\item Суточная запись сигналов разбивается на неперекрывающиеся эпохи длительностью 5~с.
	\item Эпохи, содержащие высокоамплитудные артефакты (например, вызванные движением животного), исключаются из дальнейшего анализа.
	\item Для каждой оставшейся эпохи с помощью быстрого преобразования Фурье вычисляются спектры мощности сигналов ЭКоГ и ЭСКоГ с разрешением по частоте 0.5~Гц.
	\item \textbf{Детекция МС:} В спектре мощности ЭКоГ выделяется дельта-диапазон (2.5--5.5~Гц) \cite{Fang2010}(Fang2010). Суммарная мощность в этом диапазоне образует временной ряд, в котором методом гауссовых смесей выделяются две компоненты. Компонента с более высокой средней мощностью интерпретируется как медленный сон.
	\item \textbf{Детекция ПС:} В спектре мощности ЭСКоГ выделяется тета-диапазон (5.5--8~Гц) \cite{Vyazovskiy2005}(Vyazovskiy2005). Для каждой эпохи вычисляется отношение мощности тета-ритма к мощности дельта-ритма. В полученном ряду выделяются высокоамплитудные пики, которые маркируются как эпизоды парадоксального сна.
	\item Результаты шагов 4 и 5 объединяются: эпохи, классифицированные как МС или ПС, получают соответствующие метки; все оставшиеся эпохи считаются бодрствованием.
\end{enumerate}

Результатом работы алгоритма является гипнограмма --- последовательность меток состояний для каждой 5-секундной эпохи на протяжении всей суточной записи. Именно гипнограмма служит исходным материалом для дальнейшего анализа связи структуры сна с сейсмической активностью.



В основной части работы должны быть описаны полученные результаты. Здесь возможны:
\begin{itemize}
\item определения основных понятий, систематизация имеющихся точек зрения по изучаемому вопросу, аргументированная формулировка собственной точки зрения;
\item теоретические и экспериментальные результаты, полученные автором;
\item детальные описания методов и алгоритмов решения задач и фактов, лежащих в их основе;
\item описания разработанных программ, включая образцы экранных\linebreak форм и фрагменты кода.
\end{itemize}

Помимо изложения результатов работа должна содержать методы их получения (доказательства теорем, методика проведения экспериментов, обоснования корректности алгоритмов, и т.\,п.).

В основном тексте также допустимо указывать результаты работы программ, если они имеют небольшой объём. В противном случае эти результаты (или их часть, не вошедшая в основной текст) могут быть приведены в приложениях.


Переносы в заголовках разделов не допускаются.

Необходимо обращать внимание на написание дефисов, длинных и коротких тире в основном тексте, например:

\emph{Одним из участников Великой Отечественной войны и парада Победы на Красной площади в~1945 году был выдающийся российский учёный, чьё имя с гордостью носит наш Институт, и чью столетнюю годовщину со дня рождения мы будем торжественно отмечать 21 июня~--- Иосиф Израилевич Ворович (1920--2001).}


\subsection{Рекомендуемые объёмы работ}

Рекомендуемые объёмы работ приведены в~табл.\,\ref{stud:table:1}. В~скобках указаны рекомендуемые объемы для работ в~области педагогического образования.

\begin{table}[H]
  \centering

\caption{Рекомендуемые объёмы работ}\label{stud:table:1}

\begin{tabular}{|p{85mm}|c|}
  \hline
  \multicolumn{1}{|c|}{\bf Уровень работы} &
  \multicolumn{1}{c|} {\bf Объём (в страницах)} \\
  \hline
Курсовая работа	& 8--10 (15--20) \\
  \hline
Выпускная квалификационная работа\newline
(4 курс бакалавриата)	&  15--20 (40--50)\\
\hline
Выпускная квалификационная работа\newline
(2 год магистратуры)	&  40--50 (50--70) \\
\hline
\end{tabular}
\end{table}



\subsection{Вставка изображений и таблиц}

Все таблицы, рисунки, схемы, диаграммы и другие объекты, вставляемые в текст, должны быть пронумерованы, подписаны и выровнены по центру, в~тексте должна присутствовать ссылка на них (рис.\,\ref{stud:fig:1}).


%\begin{figure}[H]
%  \centering
%  \includegraphics[scale=1.2]{Fig_T.png}
%  \caption{Диаграмма переходов машины Тьюринга}\label{stud:fig:1}
%\end{figure}


\subsection{Фрагменты исходного кода}

Для иллюстрации излагаемого материала в основной текст можно вставлять фрагменты исходного кода (тексты программ). Они должны быть набраны моноширинным шрифтом (например, \texttt{Courier New}), 12~пунктов, с~одинарным интервалом и выравниванием по левому краю. Если к~фрагментам кода требуется указывать ссылки, их тоже надо снабжать заголовками, как в~случае листингов~\ref{stud-lst:1}--\ref{stud-lst:3}, приведенных ниже.

Если строка кода не помещается на одну строку страницы, её следует разбивать на части в соответствии с принятым стилем форматирования кода, а не автоматически. Размер непрерывных фрагментов исходного кода не должен превышать половины страницы. Фрагменты большего размера следует помещать в приложении к работе.


\subsubsection{Примеры фрагментов кода}

Ниже приведены листинги с заголовками.

\begin{lstlisting}[language=C++, caption={C++, пример кода}, label=stud-lst:1]
#include <iostream>
int main()  // однострочный комментарий
{
  std::cout << "Привет, мир!" << std::endl;
}
\end{lstlisting}


\begin{lstlisting}[language=Python, caption={Python, пример кода}, label=stud-lst:2]
print("Привет, мир!")  # comments
\end{lstlisting}

\begin{lstlisting}[language=TeX, caption=\LaTeX, label=stud-lst:3]
% параметр language в наших листингах только для себя
\bf Привет, мир!
\end{lstlisting}



%=======================
\newpage
\section{Примеры оформления формул}

Внутритекстовая формула
$ \int\limits_a^b f(x)\,dx$ не всегда выглядит красиво даже в \LaTeXe.

Но формулы, вынесенные в отдельную строчку всегда выглядят шикарно
\[
  \lim_{d\to 0} S_n =
  \int\limits_a^b f(x)\,dx.
\]

Для того чтобы {\TeX} автоматически мог ссылаться на нумеруемые формулы нужно указывать команду \verb"\label"
\begin{equation}\label{eq:1}
  \frac{abc}{xyz}.
\end{equation}
Формула (\ref{eq:1}) была оформлена с помощью окружения \textsf{equation}.

%=======================
\newpage
\addcontentsline{toc}{section}{Заключение}
\section*{Заключение}

Заключение должно содержать информацию о проделанной работе и полученных результатах.

При написании текста работы следует иметь в виду, что её цель состоит в том, чтобы продемонстрировать квалификацию автора. Поэтому следует избегать общих и, тем более, тривиальных или нравоучительных высказываний. Мотивация выполняемой работы не должна носить слишком конкретный характер. Во время выступления на защите желательно избегать упоминаний об особенностях стандартных компонентов пользовательского интерфейса программ (<<нажимаем на правую кнопку>>, <<перетаскиваем фрагмент мышью>> и т.\,д.). Не следует комментировать задаваемые после защиты вопросы. Ответы на вопросы должны быть краткими.



%=======================
\newpage

\addcontentsline{toc}{section}{Литература}
\renewcommand{\refname}{\centering \textbf{Литература}}

\begin{thebibliography}{0}
\bibitem{stud:b0}
Yokoi S. et. al. Mouse circadian rhythm before the Kobe earthquake in 1995 //Bioelectromagnetics. – 2003. – Т. 24. – №. 4. – С. 289-291.

\bibitem{stud:b1}
Жуков М.\,Ю., Ширяева Е.\,В.
\LaTeXe: искусство набора и вёрстки текстов с~формулами.~-- Ростов н/Д : Изд-во ЮФУ, 2009.
\end{thebibliography}



\end{document}
% ----------------------------------------------------------------


\lstset{ %
language=C++,                 % выбор языка для подсветки (здесь это С++)
basicstyle=\small\sffamily, % размер и начертание шрифта для подсветки кода
numbers=left,               % где поставить нумерацию строк (слева\справа)
numberstyle=\tiny,           % размер шрифта для номеров строк
stepnumber=1,                   % размер шага между двумя номерами строк
numbersep=5pt,                % как далеко отстоят номера строк от подсвечиваемого кода
backgroundcolor=\color{white}, % цвет фона подсветки - используем \usepackage{color}
showspaces=false,            % показывать или нет пробелы специальными отступами
showstringspaces=false,      % показывать или нет пробелы в строках
showtabs=false,             % показывать или нет табуляцию в строках
frame=single,              % рисовать рамку вокруг кода
tabsize=2,                 % размер табуляции по умолчанию равен 2 пробелам
captionpos=t,              % позиция заголовка вверху [t] или внизу [b]
breaklines=true,           % автоматически переносить строки (да\нет)
breakatwhitespace=false, % переносить строки только если есть пробел
escapeinside={\%*}{*)}   % если нужно добавить комментарии в коде
extendedchars=true,
commentstyle=\color{mygreen},    % comment style
stringstyle=\bf,
commentstyle=\ttfamily\itshape,
keepspaces=true % пробелы между русскими буквами
aboveskip=3mm,
belowskip=3mm

}


\renewcommand\NAT@bibsetnum[1]{\settowidth\labelwidth{\@biblabel{#1}}%
   \setlength{\leftmargin}{\bibindent}\addtolength{\leftmargin}{\dimexpr\labelwidth+\labelsep\relax}%
   \setlength{\itemindent}{-\bibindent+\fivecharsapprox}%
   \setlength{\listparindent}{\itemindent}
\setlength{\itemsep}{\bibsep}\setlength{\parsep}{\z@}%
   \ifNAT@openbib
     \addtolength{\leftmargin}{\bibindent}%
     \setlength{\itemindent}{-\bibindent}%
     \setlength{\listparindent}{\itemindent}%
     \setlength{\parsep}{0pt}%
   \fi
}
\renewcommand{\thesection}{\arabic{section}.}
\renewcommand{\thesubsection}{\arabic{section}.\arabic{subsection}.}
